% Originally: content/week-03.tex, \subsection{Contraction mappings and the Banach fixed point theorem}

% Sonnet 5 (Medium)
\section{Contraction mappings and the Banach fixed point theorem}
\label{sec:banach_fixed_point}

\begin{ainote}
Not in Corsin's own notes either, but needed to justify the appeal to the \qt{Banach Fixed
Point Theorem} in the solution to \cref{ex:ai_contraction_cos} below --- and a standard companion
to completeness (\cref{def:complete_metric_space}), which is why it sits in this chapter.
\end{ainote}

\begin{definition}[Contraction mapping]
\label{def:contraction_mapping}
Let $(X,d)$ be a metric space. A map $f : X \to X$ is a \newterm{contraction}
(\germanterm{Kontraktion}) if it is Lipschitz continuous (\cref{def:lipschitz_continuous}) with
some Lipschitz constant $L \in [0,1)$, i.e.,
\[d(f(x),f(y)) \leq L \cdot d(x,y) \quad \text{for all } x,y \in X.\]
\end{definition}

\begin{theorem}[Banach fixed point theorem]
\label{thm:banach_fixed_point}
Let $(X,d)$ be a nonempty complete metric space (\cref{def:complete_metric_space}) and let
$f : X \to X$ be a contraction (\cref{def:contraction_mapping}). Then $f$ has a unique
\newterm{fixed point} (\germanterm{Fixpunkt}) $x^* \in X$, i.e., $f(x^*) = x^*$; moreover, for any
$x_0 \in X$, the iterates $x_{n+1} := f(x_n)$ converge to $x^*$.
\end{theorem}

% Supplement: Sascha Brack/Class Notes/Week_02_Notes_Monday_Updated.pdf, p. 13
\begin{proof}[Proof idea]
The core idea is to create a Cauchy sequence by repeatedly applying the contraction mapping $f$.
Fix an arbitrary starting point $x_0 \in X$, and define the sequence by $x_{n+1} := f(x_n)$ for all $n \geq 0$.
Comparing consecutive terms gives:
\[
d(x_{n+1}, x_n) = d(f(x_n), f(x_{n-1})) \leq \lambda \cdot d(x_n, x_{n-1}).
\]
Iterating this inequality yields $d(x_{n+1}, x_n) \leq \lambda^n d(x_1, x_0)$. Since $\lambda < 1$, one can use the triangle inequality and the geometric series to show that the distance between any two terms $x_m$ and $x_n$ (for $m > n$) becomes arbitrarily small as $n \to \infty$. Thus, the sequence is Cauchy.
Because $X$ is a complete metric space, this sequence converges to some limit $x^* \in X$. Since contraction mappings are continuous, one can pass the limit inside the function to conclude $f(x^*) = x^*$.
\end{proof}

% Generator: Claude Opus 5 (High) --- Sascha's proof idea establishes existence only; uniqueness
% is two lines and is half of what the theorem claims.
\begin{proof}[Uniqueness of the fixed point]
Suppose $f(x) = x$ and $f(y) = y$. Then
\[d(x,y) = d\big(f(x),f(y)\big) \leq L\,d(x,y),\]
so $(1-L)\,d(x,y) \leq 0$. Since $L<1$ we have $1-L>0$, forcing $d(x,y) \leq 0$ and hence
$d(x,y)=0$, i.e.\ $x=y$.
\end{proof}

\begin{remark}[Where each hypothesis is used]
\label{rem:banach_hypotheses}
Both hypotheses of \cref{thm:banach_fixed_point} are indispensable, and each fails in a
different way.
\begin{itemize}
  \item \textbf{Completeness} is what supplies the limit. On the incomplete space
    $X := (0,1]$, the map $f(x) := x/2$ is a contraction with $L=\tfrac12$ and has no fixed
    point in $X$ --- the iterates march towards $0$, which is missing.
  \item \textbf{$L<1$ strictly} is what forces the iterates together. The map
    $f(x) := x + \tfrac{1}{1+e^{x}}$ on the complete space $\mathbb{R}$ satisfies the strict
    inequality $|f(x)-f(y)| < |x-y|$ for $x \neq y$, yet has no fixed point, since
    $f(x)>x$ everywhere. A supremum of $1$ that is never attained is still not good enough.
\end{itemize}
The uniqueness proof above shows the second point cleanly: the whole argument is the division by
$1-L$, which is exactly what $L=1$ forbids.
\end{remark}

\begin{aiexercise}[{Contraction Mapping on $[0,1]$}]
\label{ex:ai_contraction_cos}
% Generator: Gemini 3.6 Flash (Medium)
Let $f \colon [0,1] \to [0,1]$ be defined by $f(x) := \frac{1}{2}\cos(x)$.
\begin{enumerate}[label=\textbf{(\alph*)}]
  \item Show that $f$ is a contraction mapping (\cref{def:contraction_mapping}) with Lipschitz
    constant $L = \frac{1}{2}\sin(1) < 1$.
  \item Deduce, using \cref{thm:banach_fixed_point}, that $f$ has a unique fixed point
    $x^* \in [0,1]$.
\end{enumerate}
\end{aiexercise}


% --- exercise relocated here from the dissolved sheet 3 ---

% Originally: content/exercise-sheets/week-03.tex, problem 3.5
% Source: exercises/Ex3_Analysis2_eng.pdf

\begin{exercise}[3.5 --- Multiple choice (fixed points)]
\label{ex:3.5}
Select all the statements below that are true.
\begin{enumerate}[label=\textbf{(\alph*)}]
  \item If you spread a map of Zurich on your desk, then one point of the desk will coincide with
    its representation on the map (in an ideal world).
  \item If $f : [0,1] \to [0,2]$ is $\tfrac{1}{2}$-Lipschitz, then $f$ has a fixed point.
  \item If $f : [0,2] \to [0,1]$ is $\tfrac{1}{2}$-Lipschitz, then $f$ has a fixed point.
  \item If $f : [0,1] \cup [2,3] \to [0,1] \cup [2,3]$ is continuously differentiable with
    $|f'(x)| < 1$ for all $x \in [0,1]\cup[2,3]$, then it has a fixed point.
  \item \textbf{(*)} If $f : [0,1] \to [0,1]$ is differentiable with $|f'(x)| < 1$ for all
    $x \in (0,1)$, then it has a fixed point.
\end{enumerate}
\end{exercise}
